\section*{Capítulo \ref{ch:g4:series-circuits} --- Circuitos en serie}

\begin{solution}{exo:g4:series-circuits:1}
Los jugadores se sientan en un único bucle, uno detrás de otro y sin
bifurcaciones --- como las cuentas ensartadas en un collar.
\end{solution}

\begin{solution}{exo:g4:series-circuits:2}
Todo o nada (un hueco para a todos); el orden no importa
(intercambiar jugadores no cambia nada); reparto (cuantas más
\omterm{def:g3:first-electric-circuit:bulb}{bombillas} hay en el bucle, más débil brilla cada una).
\end{solution}

\begin{solution}{exo:g4:series-circuits:3}
La \omterm{def:g3:first-electric-circuit:bulb}{bombilla} 1 se apaga también. El casquillo vacío es un hueco en el
bucle único, y la regla uno --- todo o nada --- para a todos los
jugadores a la vez.
\end{solution}

\begin{solution}{exo:g4:series-circuits:4}
No cambia nada: el \omterm{ex:g3:first-electric-circuit:switch}{interruptor} sigue mandando sobre las dos
\omterm{def:g3:first-electric-circuit:bulb}{bombillas} exactamente igual que antes. La regla dos --- el orden en
el bucle no importa.
\end{solution}

\begin{solution}{exo:g4:series-circuits:5}
Dos \omterm{def:g3:first-electric-circuit:bulb}{bombillas}: cada una claramente más floja de lo que estaba la
\omterm{def:g3:first-electric-circuit:bulb}{bombilla} sola. Tres: más flojas todavía. Lo que se reparte es el
empuje de la \omterm{def:g3:first-electric-circuit:battery}{pila} a lo largo de todo el collar.
\end{solution}

\begin{solution}{exo:g4:series-circuits:6}
Nariz contra cola: el $+$ de una tocando el $-$ de la siguiente, para
que las dos empujen en el mismo sentido a lo largo del bucle. Dibujo:
$-\,[{+}\;{-}]\,[{+}\;{-}]\,+$ en una sola fila.
\end{solution}

\begin{solution}{exo:g4:series-circuits:7}
Desde el $+$: por la tira metálica de la carcasa hasta el
pulsador, por el \omterm{ex:g3:first-electric-circuit:switch}{interruptor} hasta la \omterm{def:g3:first-electric-circuit:bulb}{bombilla}, por el hilo
incandescente de la \omterm{def:g3:first-electric-circuit:bulb}{bombilla} y de vuelta por el tubo --- a través de
la segunda \omterm{def:g3:first-electric-circuit:battery}{pila}, nariz contra cola --- hasta el $-$ de la primera
\omterm{def:g3:first-electric-circuit:battery}{pila}. Un bucle, sin bifurcaciones.
\end{solution}

\begin{solution}{exo:g4:series-circuits:8}
En cualquier punto del mismo bucle --- por ejemplo, justo después del
primer pulsador. \omterm{def:g4:series-circuits:series}{En serie}, los dos pulsadores tienen que estar
cerrados para completar la única carretera, así que el timbre suena
solo cuando se aprietan los dos a la vez.
\end{solution}

\begin{solution}{exo:g4:series-circuits:9}
Se sustituye la primera \omterm{def:g3:first-electric-circuit:bulb}{bombilla} por la que sabemos que funciona; si la
guirnalda se enciende, la culpable era la que se ha quitado. Si no,
se vuelve a poner la original y se pasa al segundo casquillo, y así
sucesivamente a lo largo de la guirnalda. En algún casquillo la
guirnalda se encenderá --- la \omterm{def:g3:first-electric-circuit:bulb}{bombilla} que se acaba de quitar es la
cuenta rota.
\end{solution}

\begin{solution}{exo:g4:series-circuits:10}
A oscuras, Tom metió una \omterm{def:g3:first-electric-circuit:battery}{pila} del revés. Nariz contra nariz, los
empujes de las dos \omterm{def:g3:first-electric-circuit:battery}{pilas} se pelean en vez de hacer cola, y la
\omterm{def:g3:first-electric-circuit:bulb}{bombilla} no recibe nada --- unas tablas a toda potencia siguen siendo
tablas. Girando una \omterm{def:g3:first-electric-circuit:battery}{pila} se recupera la cola nariz contra cola y la
\omterm{def:g1:light-and-shadows:source}{luz}.
\end{solution}

\begin{solution}{exo:g4:series-circuits:11}
Tiene razón: en un bucle único, un \omterm{ex:g3:first-electric-circuit:switch}{interruptor} colocado en
cualquier punto manda sobre todo --- el orden no importa. Con
carreteras que se bifurcan, la corriente podría rodear su
\omterm{ex:g3:first-electric-circuit:switch}{interruptor} por otra rama y el motor podría seguir en marcha; los
circuitos con ramas exigen colocar los \omterm{ex:g3:first-electric-circuit:switch}{interruptores} con más
cuidado (son los circuitos en paralelo del año que viene).
\end{solution}
